
\begin{obeassessment}
\textbf{Kuis Bab 9}

1. Jelaskan mekanisme kerja fungsi $f$ pada algoritma GOST.

2. Sebutkan parameter $w$, $r$, dan $b$ pada RC5 beserta nilai yang dibolehkan.

3. Bandingkan Serpent dan Twofish dari sisi jumlah putaran, kotak-S, dan blok
   data.

4. Hitunglah $t$, $c$, dan $n$ untuk RC5-32/12/16. Jelaskan mengapa pada kasus
   ini $n = 3t$.

5. Sebuah sistem memakai satu kunci GOST untuk mengenkripsi 100 GiB data.
   Jelaskan risiko yang timbul dan kapan kunci itu seharusnya diganti.

6. Apa perbedaan struktur antara MARS dan Serpent, dan bagaimana perbedaan itu
   tercermin pada pembagian putarannya?
\end{obeassessment}
\begin{competencychecklist}
\begin{tabularx}{\linewidth}{|X|c|}
\hline
Indikator Kompetensi & Tercapai \\
\hline
Saya dapat menjelaskan mekanisme kerja GOST dan perbedaannya dengan DES & [ ] \\
\hline
Saya dapat menjelaskan hasil kriptanalisis GOST dan batas rekeying $2^{n/2}$ blok & [ ] \\
\hline
Saya dapat menguraikan parameter fleksibel pada RC5 dan RC6 & [ ] \\
\hline
Saya dapat menjalankan jadwal kunci RC5 dan satu putaran GOST & [ ] \\
\hline
Saya dapat membedakan karakteristik Blowfish, Twofish, Serpent, dan MARS & [ ] \\
\hline
Saya dapat membedakan struktur Jaringan Feistel dan substitusi--permutasi & [ ] \\
\hline
\end{tabularx}
\end{competencychecklist}
